\section{Open Points for Confirmation}
\label{sec:openpoints}

The methodology above commits to a decision on every question that was still
open after the meeting of 4~September~2026 and the written input of
5~September~2026, so that work can proceed without waiting. Each decision and
its rationale is listed here for the supervisor to confirm or overturn. Where a
decision is overturned, the affected part of Section~\ref{sec:method} is the
only thing that changes.

\begin{enumerate}[leftmargin=*,itemsep=5pt]
  \item \textbf{Scope: document and interface analysis, not a user survey.}
    The written input of 5~September lists only document and interface sources,
    which settles the question raised in the meeting notes. The scenario-based
    online survey from the original thesis offer is treated as a possible second
    phase and is out of scope here (Section~\ref{sec:problem}).
    \emph{Confirm that this closes the question.}

  \item \textbf{Harm taxonomy: Citron and Solove as the backbone, Knabenhans et
    al.\ as the pathway annotation.} The seven categories give comparable,
    citable buckets; the harm-tree stages record whether a provider is
    describing a mechanism, a capability, or an outcome. The two are used
    together rather than merged, so that neither loses its original meaning.
    \emph{Confirm this is the intended reading of the two papers.}

  \item \textbf{\enquote{Highest rated} operationalised as: at least 1{,}000
    ratings, then ranked by average rating.} Without a minimum rating count,
    \enquote{highest rated} selects obscure apps with a handful of five-star
    reviews. \emph{Confirm the threshold, or name a different one.}

  \item \textbf{Sample size: 50 apps, stratified 30 fitness\,/\,20 fertility, as
    the final target rather than a pilot.} A separate five-app pilot is run
    first to test the instruments. \emph{Confirm 50 is the target and not a
    first pass towards something larger.}

  \item \textbf{Walkthroughs on 12 of the 50 apps, balanced six fitness\,/\,six
    fertility and three Android\,/\,three iOS per domain.} Twelve is what one
    analyst can do consistently inside the schedule; the balance keeps the
    domain and platform comparisons in RQ4 from being confounded with sample
    size. \emph{Confirm the subsample size and the balanced split, or state that
    every sampled app must be walked through --- which would require reducing
    the document sample.}

  \item \textbf{Advertisements and privacy blogs as a secondary corpus.} Both
    are coded where found, but absence is not treated as a finding, because ad
    delivery cannot be reproduced systematically.
    \emph{Confirm that they need not be in the core corpus.}

  \item \textbf{Language: German storefront, German policy where offered,
    English otherwise, with divergences recorded.}
    \emph{Confirm this is acceptable, and whether quoted material should be
    translated in the thesis or left in the original with a translation in a
    footnote.}

  \item \textbf{A second coder.} The validation design would be materially
    stronger with a second coder on part of the subsample. \emph{Can one be
    arranged through the chair or MPI-SP?} If not, the limitation is reported as
    such (Section~\ref{sec:validation}).

  \item \textbf{Choice of model for the LLM-assisted coding.} The pipeline
    requires a pinned model version that will remain available for the duration
    of the thesis and that may be used on this material. \emph{Confirm whether a
    particular model or an institutional API should be used, or whether the
    choice is the author's.}

  \item \textbf{Ethics.} No human subjects are involved, so no review is
    expected to be needed. \emph{A written confirmation from the supervisor and
    the relevant RUB\,/\,MPI-SP contact is requested before collection starts},
    to be filed in \pathname{04_Study/Ethics/}.
\end{enumerate}
